This appendix contains supplementary material for Section~\ref{sec:computing}.
We describe two algorithms in detail: step-wise AD, and our own algorithm
\replay{}. Refer to Section~\ref{sec:computing} for the notation used in this
appendix.

\subsection{Warmup: Step-wise AD}
\label{sec:ad}
We fully describe step-wise AD in Algorithm~\ref{alg:factored}. The algorithm
requires storing all $T$ optimizer states, but requires constant memory overhead
for each AD call (as each AD call is over a single step), making it feasible to
compute for small setups.

\begin{algorithm}
  \SetKwFor{For}{\textcolor{forcolor}{for}}{\textcolor{forcolor}{do}}{}
  \SetKw{Return}{\textcolor{forcolor}{Return}}{}
  \SetKwFunction{traversereverse}{\textcolor{fncolor}{reverse\_inorder\_traversal}}
  \DontPrintSemicolon
  \hspace{0pt} \tcp{Store each optimizer state on disk}
  $\{s_i\}_{i=0}^T \gets$ Train model via $A(z)$ \;
  \;
  \hspace{0pt} \tcp{Variables; shorthand for $\smash{\frac{\partial f(z)}{\partial z}}$ and $\smash{\frac{\partial f(z)}{\partial s_T}}$}
  $\bar{z} \gets 0$ \;
  $\bar{s}_T \gets \smash{\pfrac{g(s_T)}{s_T}}$ \qquad \tcp{One reverse-mode AD call} %
  \;
  \hspace{0pt} \tcp{Reverse-mode differentiate step-by-step}
  \For{$s_i \gets s_{T-1}$ \color{forcolor}{\textbf{to}} \color{black}$s_0$}{

    \hspace{0pt} \tcp{One reverse-mode AD call. Left: $\nabla_{s_i} f$. Right: contribution to $\nabla_{z} f$ at $i$.}
        $\bar{s}_i \gets \bar{s}_{i+1} \cdot \pfrac{h_{i}(s_i, z)}{s_i}, \qquad \bar{z}_i \gets \bar{s}_{i+1} \cdot \pfrac{h_{i}(s_i, z)}{z}$\;
        \;
        $\bar{z} \leftarrow \bar{z} + \bar{z}_i$ \qquad \tcp{Accumulate metagradient}
  }
  \;
  \Return $\bar{z}$ \;
  \caption{metagradients in $\mathcal{O}(T)$ space.}
  \label{alg:factored}
\end{algorithm}

\subsection{\replay{}}
\label{app:replay}
We now describe \replay{}, our method for calculating metagradients. For a free
parameter $k \in \mathbb{N}$, \replay{} requires storing $\mathcal{O}(k\log_k(T))$ optimizer
states and an additional $\mathcal{O}(\log_k(T))$ factor of computation. The
free parameter $k$ controls the trade-off between storage and required compute.
We fully describe \replay{} in Algorithm~\ref{alg:replay}. \replay{} modifies
Algorithm~\ref{alg:factored} by retrieving the optimizer states in reverse order
using a $k$-ary tree structure in lieu of a list of all the stored states.

\subsubsection{Lazy $k$-ary tree}
We now describe the $k$-ary tree structure that underlies \replay{}; for a
visual reference of this tree with $k=2$, see
Figure~\ref{fig:efficient_data_structure}. For ease of analysis we parameterize
the total number of states as $n=T+1$ (and therefore take $n-1$ total training steps) when
describing this data structure, and assume WLOG that $n$ is an integer power of
$k$. At a high level, traversing this tree recursively replays retraining to
recover all the optimizer states in reverse order, while deleting states that
are no longer needed. We call this tree ``lazy'' because it retrains only when
required to obtain states that are not yet retrieved.

The tree is a complete $k$-ary tree with $n$ leaves (and therefore $\log_k(n)$
depth) structured as follows. We start at the root, then recursively define the
rest of the tree. Every node in the tree represents a single optimizer state.
The root represents state $s_0$. To recursively define the remaining nodes: each
non-leaf node $s_i$ at depth $d$ has $k$ equally spaced (in terms of state
number) children starting---from left to right---at state $s_i$ and ending at
$s_{i + n/k^{d+1}}$. This means that the leaves correspond---from left to
right---to the states $s_0, s_1, \ldots, s_{n-1}$.

We reduce the problem of iterating over the states in reverse to the problem of
reverse in-order traversing this tree and yielding \textit{just} the
leaves---this is exactly the states in reverse order. A reverse in-order
traversal for this $k$-ary tree requires repeatedly: recursively traversing
child nodes from largest to smallest, then visiting the parent node. We design
the specifics of this traversal to maximize space and compute efficiency. To
access the children of a parent node at traversal time, we replay model training
from the smallest child state (which is stored in the parent state) to the
largest child state and store all the children. We perform this operation
recursively each time we traverse a node. After traversing the node's left side
(i.e., after ascending from this node), we delete all its child states.

Reverse in-order traversing this tree requires storing at most $k\log_k(n)$
optimizer states at a time, and in aggregate requires retraining the model
$\log_k(n)$ times. The argument for each is straightforward. Storage: the
traversal requires storing at most $k$ states for each level that it descends
(we store $k$ states whenever we first traverse to a parent node) and we remove
$k$ states for each level that the traversal ascends (we remove $k$ states after
we are done with the left traversal of a parent). Compute: we replay training to
reinstantiate the children of every parent node a single time. The $k^d$ parent
nodes at level $d$ each require replaying $\mathcal{O}(\nicefrac{n}{k^d})$
states to reinstantiate children. Therefore, in a traversal, each level requires
$\mathcal{O}(n)$ ($k^d \cdot \nicefrac{n}{k^d}$) optimizer steps. There are
$\log_k(n)$ levels with parent nodes, which means a total of $\mathcal{O}(n
\log_k(n))$ optimizer steps, or a multiplicative factor of
$\mathcal{O}(\log_k(n))$ steps compared to model training.

\begin{algorithm}
  \SetKwFor{For}{\textcolor{forcolor}{for}}{\textcolor{forcolor}{do}}{}
  \SetKw{Return}{\textcolor{forcolor}{Return}}{}
  \SetKwFunction{traversereverse}{\textcolor{fncolor}{reverse\_inorder\_traversal}}
  \DontPrintSemicolon $T \gets $ Lazy $k$-ary tree for $\mathcal{A}(z)$ \qquad \tcp{Make
  lazy $k$-ary tree of Appendix~\ref{app:replay}} \;

  \hspace{0pt} \tcp{Variables; shorthand for $\smash{\frac{\partial f(z)}{\partial z}}$ and $\smash{\frac{\partial f(z)}{\partial s_T}}$}
  $\bar{z} \gets 0$ \;
  $\bar{s}_T \gets \smash{\pfrac{g(s_T)}{s_T}}$ \qquad \tcp{One reverse-mode AD call} %
  \;
  \hspace{0pt} \tcp{Reverse-mode differentiate step-by-step; \underline{traverse $T$ instead of stored states}}
  \For{$s_i \gets s_{T-1}$ \color{forcolor}{\textbf{to}} \color{black}$s_0 \in \traversereverse{T}$}{ 
  \hspace{0pt} \tcp{One reverse-mode AD call. Left: $\nabla_{s_i} f$. Right: contribution to $\nabla_{z} f$ at $i$.}
      $\bar{s}_i \gets \bar{s}_{i+1} \cdot \pfrac{h_{i}(s_i, z)}{s_i}, \qquad \bar{z}_i \gets \bar{s}_{i+1} \cdot \pfrac{h_{i}(s_i, z)}{z}$\;
      \;
      $\bar{z} \leftarrow \bar{z} + \bar{z}_i$ \qquad \tcp{Accumulate metagradient}
  }
  \;
  \Return $\bar{z}$ \;
  \caption{\textsc{Replay}. metagradients in
  $\mathcal{O}(k\log_k(T))$ space.}
  \label{alg:replay}
\end{algorithm}

